\section{Framework and the baseline fit}
\label{sec:framework}
We use the LO BO$\chi$EFT baseline fit of Ref.~\cite{Zhang:2026iyr}
(v1, Eqs.~(12), (18), (83), (104)--(106), Tables~I, V and VI), in the
isospin limit and with full OME. The first four bare states are retained,
with a sharp cutoff $\Lambda=4$~GeV. The quark--antiquark Hamiltonian is
$H=2m_b-\nabla^2/m_b+V_0+\sigma r+\gamma/r$, with $\gamma<0$.
The bare--continuum coupling is $g=(2/\sqrt{3})\sqrt{2}\,g_l$ for
$g_l=35$~MeV; $g$ is distinct from the axial coupling $g_b=0.56$.
% Canonical entries: CONVENTIONS.md "Quark--antiquark Hamiltonian",
% "Bare--continuum vertex", "Loop measure and regulator";
% GLOSSARY.md "Baseline fit", "Isospin limit", "Compact probability".

\begin{table}[ht]
\centering
\caption{The $1^{++}$ baseline-fit inputs. Rounded values follow
Ref.~\cite{Zhang:2026iyr} (v1, Eq.~(18), Tables~I and VI);
the computation uses the full-precision values in
\texttt{code/HadronDynamics.jl/data/bochiEFT/central.toml}.
The values are fixed in $1^{++}$, not fitted anew in other sectors.}
\label{tab:baseline-inputs}
\begin{tabular}{cccccccc}
\toprule
$m_b$ (GeV) & $\sigma$ (GeV$^2$) & $\gamma$ & $V_0$ (MeV) &
$C$ (GeV$^{-2}$) & $g$ (MeV) & $N$ & $\Lambda$ (GeV) \\
\midrule
4.534 & 0.198 & $-0.434$ & 371 & $-16.27$ & 57 & 4 & 4 \\
\bottomrule
\end{tabular}
\end{table}

\begin{table}[ht]
\centering
\caption{Baseline-fit bound states from \texttt{checks/baseline\_fit.jl}.
Binding energies are relative to $m_B+m_{B^*}=10.60427$~GeV;
$\mathcal P_{\bar bb}$ is the compact probability.
The three masses associated with the $1P$- and $2P$-dominated states and
$X_b$ enter the fit, whereas the $3P$-dominated state is a prediction,
not a fourth fitted mass. The comparison is to Ref.~\cite{Zhang:2026iyr}
(v1, Table~V dagger rows and Table~VI, $\Lambda=4$~GeV).}
\label{tab:baseline-states}
\begin{tabular}{lrrr}
\toprule
State & $M$ (GeV) & $E_B$ (MeV) & $\mathcal P_{\bar bb}$ (\%) \\
\midrule
$\chi_{b1}(1P)$-dominated & 9.893264 & 711.006 & 99.638 \\
$\chi_{b1}(2P)$-dominated & 10.255276 & 348.994 & 98.301 \\
$X_b$ & 10.513270 & 91.000 & 1.887 \\
$\chi_{b1}(3P)$-dominated & 10.543615 & 60.655 & 85.126 \\
\bottomrule
\end{tabular}
\end{table}

The literature comparison passes for all four masses within 1~MeV and
all four component probabilities within one percentage point absolute.
The two-body $D$-wave probability is the sum of the
$B\bar B^*({}^3D_1)$ and $B^*\bar B^*({}^5D_1)$ weights, and the
three-body probability is the negative expectation of the energy derivative
of the OME potential, as required by Ref.~\cite{Zhang:2026iyr}
(v1, Eqs.~(104)--(106), Appendix~F). These model-dependent weights are
not identified with the $S$-wave compositeness $X_S$ of Eq.~(107).
At the reference resolution ($n_q=24$, Cornell grid 12000 points,
transform quadrature 800 points), doubling all three resolutions changes
the masses and binding energies by at most $0.007$~MeV and any component
probability by at most $0.001$ percentage points (rounded upward).
These shifts estimate numerical resolution error for this comparison,
not the LO theory uncertainty; the box and tail settings are unchanged.
The refinement command and individual shifts are recorded in
\texttt{plan/slices/01-checks-env-baseline-fit.md}.
